\section{Segnali e servizi multimediali}
\subsection{Segnali multimediali}
\subsubsection*{Medium e multimedia}
\begin{itemize}
	\item un medium (o media al plurale) è un mezzo per rappresentare, immagazzinare e trasportare informazioni, esempi di
	media sono testo, suono, immagini, video, ecc.
	\item un multimedia è un sistema di rappresentazione/immagazzinamento/trasporto di informazioni che utilizza più media
	diversi
\end{itemize}

\subsubsection*{Segnali e comunicazioni multimediali}
\begin{itemize}
	\item un segnale è una funzione di una o più variabili dipendenti in funzione di una o più variabili indipendenti
	(tempo per audio, spazio per immagini, spazio-tempo per video, ecc.), può essere analogico (dominio e immagine
	continui) o digitale (dominio e immagine discreti) e dal punto di vista ingegneristico è un mezzo che trasporta
	informazioni
	\item un segnale multimediale è la rappresentazione di uno o più media (audio, video, immagini, ecc.)
	come segnali digitali (discreti e quantizzati) che trasmettono informazioni
	\item una comunicazione multimediale è la trasmissione di segnali multimediali
\end{itemize}

\subsubsection*{Digitalizzazione dei segnali analogici e multimediali}
Per rappresentare delle informazioni provenienti da segnali analogici (es. suono, immagini) in segnali multimediali è
necessario eseguire le operazioni di campionamento, quantizzazione e codifica.

Il processo di digitalizzazione comporta inevitabilmente una perdita di informazione. In quelli attualmente utilizzati,
però, non è possibile percepire tale perdita. Per questo motivo i segnali multimediali digitali saranno considerati
``segnali originali'' in quanto indistinguibili da quelli analogici di partenza.

\subsubsection*{Sorgenti statiche e dinamiche}
\begin{itemize}
	\item una sorgente statica produce informazioni tempo invarianti (es. immagini, testi), supportano bene jitter, latenze
	e variazioni di bitrate, ma non supportano bene errori di trasmissione
	\item una sorgente dinamica produce informazioni in funzione del tempo (es. audio, video), richiedono sincronismo tra
	diversi media (es. audio, video e testo dei sottotitoli), supportano bene errori di trasmissione, ma sono sensibili
	a jitter e latenza	
\end{itemize}

\subsection{Servizi multimediali}
\subsubsection*{Categorie dei servizi multimediali}
Un servizio multimediale è un sistema di rappresentazione, immagazzinamento e trasporto di informazioni basato su segnali
multimediali. I servizi multimediali possono essere classificati in streaming e bulk transfer, in base a come i dati
trasmessi vengono fruiti dall'utente finale.

\subsubsection*{Streaming}
Un servizio streaming permette la fruizione dei media anche prima della completa ricezione dei dati. È composto da una
sorgente dinamica che invia i dati a passo costante (stream) a uno o più ricevitori che li memorizzano in un buffer
(playout) per poi riprodurli. Richiede una attenta gestione del buffer, della banda, della latenza, del jitter e
dell'ordine con cui i pacchetti vengono ricevuti. In base al tipo di fruizione dei media, lo streaming può essere:
\begin{itemize}
	\item \textbf{stored}: il contenuto viene trasmesso in maniera asincrona, dopo essere stato prodotto, elaborato
	e memorizzato in una memoria dedicata (server); es. video on demand come netflix e youtube, podcast e musica
	come spotify, ecc.
	\item \textbf{live}: il contenuto viene trasmesso appena viene prodotto, con latenze di poche decine di secondi;
	es. eventi sportivi in diretta, ecc.
	\item \textbf{real-time}: caso particolare dei servizi live streaming in cui la latenza non deve superare i 150 ms;
	es. videochiamate, teleconferenze, giochi online, ecc.
\end{itemize}

\noindent
I servizi di streaming in cui l'utente può interagire con il contenuto multimediale sono detti interattivi.
\begin{itemize}
	\item \textbf{streaming interattivo}: l'utente può interagire con il contenuto multimediale, ad esempio attraverso
	chat, votazioni, commenti, domande in trasmissioni live o mettere in pausa, riavvolgere o saltare parti del contenuto
	in servizi stored
	\item \textbf{streaming non interattivo}: l'utente non può interagire con il contenuto multimediale, ad esempio nei
	servizi real time
\end{itemize}
	
\subsubsection*{Bulk transfer}
Un servizio bulk transfer permette la fruizione dei media solo dopo la completa ricezione dei dati. Sono costituiti
principalmente da sorgenti statiche, per cui tollerano bene latenze e jitter, ma non tollerano errori di trasmissione.
Come per i servizi streaming, anche i servizi bulk transfer sono classificati in:
\begin{itemize}
	\item \textbf{stored}: il contenuto viene trasmesso in maniera asincrona, es. pagine web, download di file, ecc.
	\item \textbf{live}: il contenuto viene trasmesso appena viene prodotto, es. messaggistica, email, ecc.
	\item non esistono servizi bulk transfer real-time
\end{itemize}

\subsection{Requisiti e obiettivi dei servizi multimediali}
\subsubsection*{Parametri di misurazione dei requisiti}
\begin{itemize}
	\item banda/tasso/throughput: quantità di informazione trasmessa per unità di tempo, in bit/s
	\item latenza: ritardo dovuto alla trasmissione dei dati, in secondi
	\item jitter: variazione della latenza, in secondi
	\item tasso d'errore: percentuale di bit errati rispetto al totale dei bit trasmessi, in \%
	\item complessità: costo computazionale, consumo, ritardi di computazione, \dots
	\item fedeltà al segnale originale: compromesso tra qualità e tasso d'errore
\end{itemize}

\subsubsection*{Obiettivo dei servizi multimediali}
L'obiettivo dei servizi multimediali è quello di utilizzare un'infrastruttura di rete già presente (es. rete internet)
e sviluppare opportuni protocolli per privilegiare determinate esigenze (banda, latenza) rispetto ad altre, per garantire
la migliore qualità possibile in funzione dei limiti strutturali imposti dalla rete.

\subsubsection*{Classificazione banda-latenza}
\begin{center}
	\begin{minipage}{0.44\textwidth}
		\begin{itemize}
			\item bassa banda e latenza: controllo remoto
			\item bassa banda, alta latenza: IoT, sensori
		\end{itemize}
	\end{minipage}
	\hfill
	\begin{minipage}{0.55\textwidth}
		\begin{itemize}
			\item alta banda, bassa latenza: AR/VR, guida autonoma
			\item alta banda, alta latenza: streaming video
		\end{itemize}
	\end{minipage}
\end{center}

\subsubsection*{Servizi URLL - Ultra Reliable Low Latency}
Esistono alcuni servizi che richiedono latenza e tasso d'errore estremamente bassi come ad esempio:
\begin{itemize}
	\item automazione industriale, controllo del movimento (latenza \(\leq 1\) ms, tasso d'errore \(\leq 10^{-9}\))
	\item tactile internet (latenza \(\leq 1\) ms, tasso d'errore \(\leq 10^{-4}\))
	\item smart grid, controllo di sistemi energetici (latenza \(\leq 3-5\) ms, tasso d'errore \(\leq 10^{-7}\))
	\item intelligent transport system, comunicazione tra veicoli (latenza \(\leq 5\) ms, tasso d'errore \(\leq 10^{-4}\))
	\item guida autonoma (latenza \(\leq 15 - 20\) ms, tasso d'errore \(\leq 10^{-4}\))
	\item controllo remoto (latenza \(\leq 5-100\) ms, tasso d'errore \(\leq 10^{-2}\))
	\item automazione processi industriali, controllo di macchinari (latenza \(\leq 100\) ms, tasso d'errore \(\leq 10^{-5}\))
\end{itemize}
